\chapter{Disasters-Natural and Human-made}
\section{Get Rolling}
\textbf{Look at the picture an answer the questions below.}
\begin{enumerate}
    \item What do you see in the picture?
    \item What difficulties would people fact in such situation?
\end{enumerate}
\section{What to Know}
Disasters, which can be natural or human-made, can cause a lot of damage, but we can be prepared and take certain measures to reduce measures to reduce their impact.
\section{\textbf{Learning Objectives}}
\begin{itemize}
    \item to identify various types of natural and human-made disasters
    \item to differntiate between accidents and disasters
    \item to distinguish between natural and human-made disasters
    \item to identify causes and effects of natural and human-made disasters
    \item to explain the safety measures to be taken during various disasters
    \item to find our the role of emergency services and voluntary organisations during a disaster
\end{itemize}
We often hear about unexpected even like floods, earthquakes, forest fires, chemical spills or crumbling of buildings happening around the world. Such unexpected events usually kill a lot of people and cause a lot of damage. Hence, any unexpected event that causes a massive destruction to life and property and often also to the environment is called a \textbf{disaster.}
\vspace{0.1in}
\par
Sudden events that cause injury to people or damage to property are called \textbf{accidents.} An accident can happen anywhere if caution slips. It can happen due to human carelessness or equipment failure. For example, a fall during trekking, an engine failure or vehicle tyre burst. Small things can also cause accidents. For example, a minor accident can occur if a water bottle falls on someone's foot.
\vspace{3.0in}
\section{Disaster}
Disasters are of two types-natural and human-made-depending upon their causes.
\subsection{Natural Disaster}
Nature has an equal power of creation and destruction. A \textbf{natural disaster} is caused by the forces of nature, such as heavy rains and strong winds. Floods, earthquakes, tsunamis, droughts, cyclones, volcanic eruptions are some examples of natural disasters. Human beings cannot control natural disasters.
\vspace{0.1in}
\par
Human-made disasters are caused by the actions of humans. Rail accidents, air crashes, chemical spills, wars, etc. are human-made disasters.
\subsection{\textbf{Earthquakes}}
\textbf{Earthquakes} are one of the most destructive natural disasters that can strike a place. They cause massive destruction of life, property and the natural environment.
\vspace{0.1in}
\par
An earthquake is caused by sudden movements under Earth's surface that cause everything on Earth to move. It starts at a point within Earth, which is called the \textbf{focus.} The point exactly above the focus on Earth's surface is called the \textbf{epicentre.} The epicentre suffers most damage during a earthquake. A \textbf{seismograph} is an instrument used to record an earthquake tremors. The scientific study of earthquakes is called \textbf{sesimology.}
\vspace{0.4in}
\par
The 2015 earthquake in Nepal caused huge destruction of life and property. More than 8,000 people were killed and thousands became homeless.
\vspace{0.1in}
\par
Even mild earthquakes can cause heavy destruction and affect the mental and physical health of people. We cannot stop earthquakes. However, we can take some measures to save ourselves from such disasters.
\vspace{0.1in}
\par
Safety measures to be followed during the earthquakes:
\begin{itemize}
    \item Immediately move to open area away from trees, electric and telephone poles and buildings.
    \item Use the stairs and not elevators during and immediately after an earthquake. Elevators may stop as power may go off during an earthquake.
    \item If you are inside a building, drop to the floor, take cover the under a strong piece of furniture like a desk, table or bed and hold on to it tightly.
    \vspace{1.6in}
    If you are in a car, stop it away from trees, buildings and hoardings, and wait inside it.
    \item Never use matchsticks, gas stoves or light switches during an earthquake, as a gas leak or broken electric wires can start a fire.
\end{itemize}
If there is a fire in a building during an earthquake, follow these measures for your safety:
\begin{itemize}
    \item Leave the building immediately.
    \item Do not use elevator.
    \vspace{0.4in}
    \item If you see smoke under a door, do not go out of that door. 
    \item Before opening a door, touch it with the back of your hard. Don't open the door if it feels hot.
    \item If there is smoke, crawl low underit. Cover your head, mouth and nose with a wet cloth so that smoke does not enter your lungs. Breathing is smoke is harmful.
    \item Dial 101 to call the fire department. You can dial 112 in all kinds of emergencies.
    \item If ;your clothes catch fire, don't run. Running will flare up the fire. Get down on the floor and roll over. This will help to put out the flames.
\end{itemize}
\vspace{1.8in}
\subsection{\textbf{Volcanic eruptions}}
Deep inside Earth are some hot gases and molten rocks, which come out at times through an opening at the top of a mountain or a hill. Such mountains or hills are known as \textbf{volcanoes.}
The coming out of hot gases, ash and molten rocks from volcanoes is called a volcanic eruption.
\vspace{0.1in}
\par
Some volcanoes erupt very often, some erupt once in many years, and some are not likely to erupt again.
\vspace{0.1in}
\par
Volcanic eruptions can cause heavy loss of life and property and air pollution. They dangerous in densely populated regions of the world.
\vspace{0.5in}
\par
Most of the time, scientists can predict a volcanic eruption. So, it is possible to save ourselves from volcanic eruptions most of the time. But sometimes, volcanoes may erupt without any warning. So, people who live near a volcano and tourists visiting an area with a a volcano need to follow some safety measures before and during a volcanic eruption.
\vspace{0.1in}
\par
I wonder where the largest active volcano on Earth is located!
\vspace{0.1in}
\par
Safety measures to be followed during and after volcanic eruptions:
\begin{itemize}
    \item Follow the guidance given on radio or television.
    \item Seal all doors and windows of your house or office so that hot air, gases and ash do not enter the building.
    \item Cover your nose with a piece of wet cloth so that you do not inhale smoke.
    \item Avoid low-lying areas where molten lava can flow.
    \item After the eruption, get medical treatment for any injuries, breathing problems or irritation in the ears, eyes and nose.
\end{itemize}
\subsection{\textbf{Tsunamis}}
The word 'tsunami' has its origin in the Japanese words \textit{tsu} meaning harbour and \textit{nami} meaning wave. A tsunami is a series of fast moving-massive waves.
\vspace{0.1in}
\par
A tsunami is mostly caused by volcanic eruptions and earthquakes under the sea. These waves are not very high in the ocean. But when the waves travel inland, they increase their height dramatically and cause heavy destruction to life and property.
\vspace{0.1in}
\par
A tsunami can happen any time during the day or night and hence is dangerous for the fishing community, people living in the coastal areas and tourists visiting their places. We need to keep a few things in mind while living or visiting a tsunami-prone area.
\vspace{0.1in}
\par
Safety measures to be taken during a tsunami:
\begin{itemize}
    \item The speed of a tsunami is the same as that of a bullet train. So, if you feel a strong earthquake, move quickly to higher ground or land (away from the sea) immediately.
    \item Just before a tsunami, the water goes back to the sea and the sea floor can be seen. In such an event, leave the shore immediately and rush to higher ground.
    \vspace{0.4in}
    \item Don't ignore the strange behaviour of animals that may be a warning of an upcoming tsunami.
    \item Remain on higher ground! Tsunamis last for hours.
    \item Hold on to anything to protect yourself from being washed away by the waves.
\end{itemize}
\subsection{\textbf{Floods}}
Floods are the most common weather-related natural disasters. Floods occur when rivers rise above their banks and the water overflows into nearby land. Floods generally occur when there is very heavy rainfall during the rainy season or monsoons. Floods are also caused by hurricanes, tsunamis, storms, broken dams, cloud bursts and the sudden melting of snow on mountains during the summer season. Sometimes, floods happen very suddenly and with little warning. These are called \textbf{flash floods.} They are very dangerous. Floods damage human and animal lives, crops, roads, houses and other property.
\vspace{0.1in}
\par
In June 2013 thousands of locals and pilgrims visiting Kedarnath Temple died in the flash floods in Uttarakhand.
\vspace{0.1in}
\par
We can stay safe during floods if we follow some safety measures.
\vspace{0.1in}
\par
Safety measures to be followed during floods:
\begin{itemize}
    \item Stay informed about the latest news on floods in your area.
    \item Do not walk through flowing water. If you have to, use a stick to check the firmness of the ground.
    \item Don not drive a vehicle through flooded areas.
    \item Climb onto the roof if your house gets flooded.
    \item Stay away from electrical poles and wires.
    \item Get your emergency kit ready. It should have a battery-operated radio, dry food items, clean bottled water, a first a kit, a strong rope, a whistle and a cutter.
\end{itemize}
\vspace{0.4in}
\subsection{\textbf{Cyclones}}
A \textbf{cyclone} is a spinning storm that starts in the sea and moves towards the coast.
\vspace{0.1in}
\par
High-speed winds and heavy rains occur during a cyclone, which uproots trees, destroys houses and can kill people and animals. The coastal areas get flooded as these strong winds cause high waves in the sea. Cyclones hit only the warm parts of the world and do not affect the cold parts.
\vspace{0.1in}
\par
A very powerful cyclone name 'Hudhud' hit India's east coast at Vishakhapatnam in October 2014. It caused extensive damage to the coastal districts of Andhra Pradesh and Odisha.
\vspace{0.1in}
\par
Safety measures to be followed during a cyclone:
\begin{itemize}
    \item Prepare and emergency kit with food and water to last for at least three days.
    \item Listen to the radio and watch the television to stay informed of the weather conditions.
    \item Stay away from the sea or the ocean and move to safer areas before the cyclone arrives.
\end{itemize}
\subsection{Human-made Disasters}
Human-made disasters occur due to human activities. Some human activities are directly responsible for disasters. Industrial disasters like the Bhopal gas tragedy, wars, environmental pollutions (that is, those of noise, air, water and soil) are some of the disasters that result from direct human actions. Some human activities like mining, deforestation and construction in fragile areas are also directly or indirectly responsible for certain disasters such as earthquakes, landslides and floods. Both the frequency and intensity of these disasters have increased over the years.
\vspace{1.6in}
\subsection{Helping Survivors}
Mist disasters, whether natural or human-made, come without a warning. They leave people injured and homeless. Who do you think looks after the survivors of disasters? Many government agencies and non-governmental organisations work together to provide help to the victims of natural and human-made disasters. In India, thee government has set up the National Disaster Response Fund and the State Disaster Response Fund under the Disaster Management Act, 2005. The disasters covered under this scheme are cyclones, droughts, earthquakes, fires, floods, tsunamis, hailstorms, landslides, avalanches, cloudbursts, pest attacks, frost and cold waves. The primary aim of disaster response is rescuing survivors from immediate danger and stabilising the physical and emotional conditions of survivors. Emergency services such as hospitals, ambulances, fire engines, police, air force, navy and army, quickly start helping and rescuing people. The survivors are provided with food, clothing and temporary shelters. The people who are unwell injured are brought to hospitals, while the police makes sure that peace is maintained in the areas affected by the disaster. International help in the form of food, medicines and money also comes in from other countries and the \textbf{United nations.} Voluntary organisations such as the \textbf{Red Cross} and the \textbf{World Health Organisation} help in providing food and medical services to the injured. They give people care and support to start a new life again.
\vspace{2.8in}
\section{Now I Know}
\begin{itemize}
    \item A sudden, unexpected event that causes harm and injury to people or causes damage to property is called an accident.
    \item An event that causes massive destruction to life and property and also sometimes harms the environment is called a disaster.
    \item Disaster are of two types-natural disasters and human-made disasters.
    \item Disaster causes by the forces of nature such as earthquakes and volcanic eruptions are called natural disasters.
    \item Disasters caused directly or indirectly by human activities such as bomb explosions and train accidents are called human-made disasters.
    \item Earthquakes are the most destructive of natural disasters.
    \item In India, the government has set up the National Disaster Response Fund and the State Disaster Response Fund under the Disaster Management Act, 2005.
    \item After a disaster, emergency services, like hospitals, ambulances, fire engines, police, the air force, the navy and the army start helping prople immediately.
    \item We can also provide help, care and support to the survivors of disasters.
\end{itemize}